%! Linear Regression — Unit 1, Lesson 1.5 — Homework
% Gradual release: the "you do alone" block. Homework is generated per lesson (user direction,
% 2026-08-31); the teacher may substitute a DeltaMath set where the content is available there.
\documentclass[10pt]{article}
\usepackage{algebra2-article}
\usepackage{algebra2-boxes}
% Coordinate grid: {xmin}{xmax}{ymin}{ymax}
\newcommand{\lgrid}[4]{%
  \draw[step=1, linegray!40, very thin] (#1,#3) grid (#2,#4);
  \draw[->, semithick, charcoal] (#1,0)--(#2,0) node[below right, font=\scriptsize]{$x$};
  \draw[->, semithick, charcoal] (0,#3)--(0,#4) node[left, font=\scriptsize]{$y$};}
\newcommand{\dpt}[2]{\fill[navy] (#1,#2) circle (2.6pt);}

\begin{document}

\pageheader{Unit 1, Lesson 1.5}{Homework}

\vspace{0.10in}

\begin{notesbox}{Practice}
Show your work. Read direction from the \emph{sign} of $r$ and strength from its \emph{size};
give every slope its units; and check each prediction against the range of the data.

\begin{enumerate}[leftmargin=1.9em, itemsep=10pt]
  \item \textbf{Match $r$ to the plot.} Assign each value $r\in\{\,0.98,\ -0.60,\ 0.05\,\}$ to a
        plot, and name each plot's direction.
    \begin{center}
    \begin{tikzpicture}[scale=0.32]
      \lgrid{-0.5}{6}{-0.5}{7}
      \dpt{1}{1}\dpt{2}{2}\dpt{3}{3}\dpt{4}{5}\dpt{5}{5}
      \node[font=\scriptsize\bfseries] at (2.5,-1.8){(P)};
    \end{tikzpicture}
    \hspace{0.4cm}
    \begin{tikzpicture}[scale=0.32]
      \lgrid{-0.5}{6}{-0.5}{7}
      \dpt{1}{5}\dpt{2}{3}\dpt{3}{4}\dpt{4}{2}\dpt{5}{3}
      \node[font=\scriptsize\bfseries] at (2.5,-1.8){(Q)};
    \end{tikzpicture}
    \hspace{0.4cm}
    \begin{tikzpicture}[scale=0.32]
      \lgrid{-0.5}{6}{-0.5}{7}
      \dpt{1}{3}\dpt{2}{5}\dpt{3}{2}\dpt{4}{4}\dpt{5}{3}
      \node[font=\scriptsize\bfseries] at (2.5,-1.8){(R)};
    \end{tikzpicture}
    \end{center}
    (P) $r=$ \blank{1.2cm}, direction \blank{1.4cm} \quad (Q) $r=$ \blank{1.2cm}, direction
    \blank{1.4cm} \quad (R) $r=$ \blank{1.2cm}, direction \blank{1.4cm}

  \boxguard[16]
  \item \textbf{Sign versus size.} Two studies report $r=-0.85$ and $r=0.45$.
    \begin{center}
    \renewcommand{\arraystretch}{1.4}
    \begin{tabularx}{\linewidth}{@{}X|X@{}}
    \textbf{(a) $r=-0.85$} & \textbf{(b) $r=0.45$} \\
    \hline
    Direction \blank{1.6cm}; \; strength \blank{1.4cm} & Direction \blank{1.6cm}; \; strength \blank{1.4cm} \\
    \end{tabularx}
    \end{center}
    (c) A classmate says ``$r=-0.85$ is the weaker one, because it is negative.'' What is wrong
    with that? \writelines{2}

  \boxguard[24]
  \item \textbf{From a data table.} A biologist measures a seedling's height every five days.
    \begin{center}
    \renewcommand{\arraystretch}{1.2}
    \begin{tabular}{|l|c|c|c|c|c|c|c|}
    \hline
    day $x$ & $0$ & $5$ & $10$ & $15$ & $20$ & $25$ & $30$ \\ \hline
    height $y$ (cm) & $5$ & $10$ & $11$ & $18$ & $20$ & $26$ & $29$ \\ \hline
    \end{tabular}
    \end{center}
    Technology gives $\hat y=0.8x+5$ with $r=0.99$.
    \begin{enumerate}[label=(\alph*), leftmargin=1.9em, itemsep=4pt]
      \item The slope $0.8$ means: each additional \blank{1.0cm}, the height grows by about
            \blank{0.9cm} \blank{1.0cm} (units!). \; The intercept $5$ is the predicted height on
            day \blank{0.8cm}.
      \item Predict the height on day $20$, then compare with the table.
        \begin{work}
          \hat y &= 0.8(20) + 5 \\
                 &= 16 + 5 \\
                 &= 21
        \end{work}
        The table says \blank{0.9cm} cm; the prediction is \blank{1.2cm} (interpolation /
        extrapolation), and it is close.
      \item Predict the height on day $100$.
        \begin{work}
          \hat y &= 0.8(100) + 5 \\
                 &= 80 + 5 \\
                 &= 85
        \end{work}
        This is \blank{2.2cm}. Why should the biologist not trust it? \blank{6.0cm}
    \end{enumerate}

  \boxguard[16]
  \item \textbf{The special values.} Say what the scatterplot looks like, or why it cannot happen.
    \begin{enumerate}[label=(\alph*), leftmargin=1.9em, itemsep=4pt]
      \item $r=1$ exactly: the points \blank{5.0cm}
      \item $r=0$: does this prove the two variables are unrelated? \blank{0.9cm} \; A tight
            \emph{curved} pattern can have $r\approx0$, because $r$ measures only a
            \blank{1.4cm} fit.
      \item A study reports $r=1.3$. What must be true? \blank{4.6cm}
    \end{enumerate}

  \boxguard[18]
  \item \textbf{Correlation versus causation.} Across an elementary school, \emph{shoe size} and
        \emph{reading level} have a strong positive correlation.
    \begin{enumerate}[label=(\alph*), leftmargin=1.9em, itemsep=4pt]
      \item Does a bigger shoe \emph{cause} better reading? \blank{0.9cm} \; Name the lurking
            variable that pushes both up: \blank{2.8cm}
      \item If the study looked only at $9$-year-olds, would you expect the correlation to stay
            strong? Explain. \writelines{2}
    \end{enumerate}

  \item \textbf{Multiple choice (SOL-style).} Which correlation coefficient indicates the
        \emph{strongest} linear relationship?
    \begin{enumerate}[label=\Alph*., leftmargin=2.2em, itemsep=1pt]
      \item $r=0.35$
      \item $r=-0.82$
      \item $r=0.60$
      \item $r=-0.10$
    \end{enumerate}
    Say in one sentence how you ruled out one of the others. \writelines{2}
\end{enumerate}
\end{notesbox}

\boxguard[20]
\begin{extensionbox}
\textbf{Which line fits better?} For the five data points below, two candidate lines are proposed:
$y_1=1.5x+0.5$ and $y_2=x+2$. Fill in each line's predictions and its \emph{error} (error $=$
data $-$ prediction) at every point.
\begin{center}
\renewcommand{\arraystretch}{1.25}
\begin{tabular}{|l|c|c|c|c|c|}
\hline
$x$ & $1$ & $2$ & $3$ & $4$ & $5$ \\ \hline
data $y$ & $2$ & $3$ & $5$ & $6$ & $8$ \\ \hline
$y_1=1.5x+0.5$ & \blank{0.8cm} & \blank{0.8cm} & \blank{0.8cm} & \blank{0.8cm} & \blank{0.8cm} \\ \hline
error & \blank{0.8cm} & \blank{0.8cm} & \blank{0.8cm} & \blank{0.8cm} & \blank{0.8cm} \\ \hline
$y_2=x+2$ & \blank{0.8cm} & \blank{0.8cm} & \blank{0.8cm} & \blank{0.8cm} & \blank{0.8cm} \\ \hline
error & \blank{0.8cm} & \blank{0.8cm} & \blank{0.8cm} & \blank{0.8cm} & \blank{0.8cm} \\ \hline
\end{tabular}
\end{center}
Total error \emph{size} (add the absolute values): \; $y_1$: \blank{1.2cm} \quad $y_2$:
\blank{1.2cm}. \; The better fit is \blank{0.9cm}, because its predictions are \blank{2.4cm}
to the data overall --- which is exactly what technology minimises.
\end{extensionbox}

\begin{spiralbox}
\textbf{Coming next --- the Unit 1 Test, then Unit 2: Quadratic Functions.} You have now met
linear functions from every angle: characteristics, forms, absolute value, piecewise, and fitting
a line to data. In Unit~2 the data can \emph{curve} --- a falling-then-rising pattern no line can
capture --- and the curve of best fit becomes a \textbf{parabola}.
\end{spiralbox}

\end{document}
